\begin{abstract}
Coarse-to-fine next-concept models (NCM) generate text by emitting one level of a multiscale residual vector quantizer (MS-RVQ) per parallel step. Two training-free extensions seem natural for such models: predictor-corrector revision of the cumulative latent, borrowed from score-based diffusion samplers, and adaptive early halting, motivated by the nested fine-drop training that makes intermediate latents decodable. We ask whether either can improve a frozen NCM. We first derive the conditions under which each mechanism can succeed: a corrector reduces accumulated drift only if it is contractive \emph{in closed loop} on the generator's support (Proposition~1), and halting at level $k^\ast$ bounds the output KL divergence by $L_D^2\tau^2/(2(1-\rho)^2)$ only if latent increments decay geometrically with $\rho<1$ and the final decoding has margin (Proposition~2). We then test these conditions on five official NCM checkpoints (Sudoku easy/medium/hard, WMT14 De--En, XSum) with a bit-exact harness, five seeds and paired tests, and identify three properties of standard MS-RVQ pre-training that violate them. (i) Off-manifold exposure bias: continuous correction lowers open-loop latent error on WMT14 (\WmtMseLo--\WmtMseHi{} vs.\ \WmtUncLo--\WmtUncHi) yet collapses generation (Sudoku-hard \LoopVoneHard\%, WMT14 \LoopVoneWmt{} BLEU, XSum \LoopVoneXsum{} ROUGE-L); codebook projection removes most of the collapse but does not beat NCM, and DAgger fine-tuning of the generator gives no significant gain (and a significant loss on Sudoku-easy, where validation is saturated). (ii) Non-decaying flux: $\hat\rho>1$ on \RhoGtOneSud{}/10 Sudoku-hard and \RhoGtOneWmt{}/9 WMT14 transitions. (iii) Terminal-level concentration: no Sudoku board's intermediate decoding equals its final one, so even per-sequence margin halting cannot exit early. We conclude with pre-training prescriptions that would satisfy both propositions.
\end{abstract}
